\chapter{Explainability della politica appresa}
\label{cap:xai}

\section{Motivazione}
\label{sec:xai_motivazione}

Il Capitolo~\ref{cap:rl} stabilisce \emph{che} la politica funziona: sul protocollo
held-out il modello a doppia corona con $v_{\max}=2\,$m/s risolve il 99{,}7\% degli
episodi. Non dice per\`o \emph{perch\'e} funzioni, n\'e su quale parte delle 196
dimensioni dell'osservazione si regga. \`E una lacuna che conta in un contesto
operativo: sapere quali grandezze la politica sfrutta davvero significa sapere quali
sensori sono indispensabili a bordo del veicolo, quali sono ridondanti e quali
degradazioni sensoristiche la missione tollera.

Il capitolo separa due domande che la letteratura tende a confondere:

\begin{itemize}[leftmargin=1.4em, itemsep=2pt]
  \item \textbf{Che cosa la rete guarda} --- quali blocchi dell'osservazione
        influenzano l'uscita della politica. \`E una domanda \emph{osservazionale}
        sulla funzione appresa, e si risponde con l'attribuzione
        (Sezione~\ref{sec:xai_shapley}).
  \item \textbf{Che cosa serve al compito} --- quali blocchi, se rimossi, degradano
        la prestazione. \`E una domanda \emph{causale} sul task, e si risponde con
        un intervento (Sezione~\ref{sec:xai_ablation}).
\end{itemize}

Le due domande non hanno la stessa risposta, e il punto interessante \`e proprio
dove divergono: una feature pu\`o avere attribuzione alta e impatto nullo perch\'e
ridondante rispetto ad altre, oppure attribuzione modesta ed effetto grande perch\'e
non sostituibile da nessun'altra. Il capitolo si chiude
(Sezione~\ref{sec:xai_sintesi}) mettendo a confronto le due letture.

\section{Metodi post-hoc: rassegna e criterio di scelta}
\label{sec:xai_rassegna}

Lo strumentario standard dell'explainability post-hoc model-agnostic \`e essenzialmente
duplice. \textbf{LIME}~\cite{lime} approssima il modello nell'intorno di una singola
predizione con un surrogato lineare sparso, stimato su perturbazioni del punto pesate
da un kernel di prossimit\`a. \textbf{SHAP}~\cite{shap} riconduce invece un'intera
famiglia di metodi al \emph{valore di Shapley}~\cite{shapley1953} della teoria dei
giochi cooperativi: le feature sono i giocatori, la predizione \`e il payoff da
ripartire, e il valore di Shapley \`e l'unica ripartizione che soddisfa gli assiomi di
efficienza, simmetria, dummy e additivit\`a. Accanto a questi si colloca la famiglia
delle attribuzioni basate sul gradiente, di cui gli \emph{Integrated
Gradients}~\cite{integrated_gradients} sono il rappresentante assiomatico.

\`E importante distinguere due livelli che il nome ``SHAP'' tende a sovrapporre: il
valore di Shapley \`e il \emph{concetto}, mentre KernelSHAP, DeepSHAP e TreeSHAP sono
\emph{stimatori} di quel valore, resi necessari dal fatto che il numero di coalizioni
cresce come $2^n$ nel numero $n$ di giocatori. La distinzione \`e decisiva per il
seguito.

Nel contesto di questa tesi l'applicazione meccanica di questi strumenti sarebbe
inadeguata, per tre ragioni.

\paragraph{Granularit\`a e collinearit\`a.} L'osservazione non \`e un insieme di 196
grandezze indipendenti ma un insieme di \emph{blocchi} fortemente collineari al loro
interno: 72 valori sono la storia della corona interna, altri 72 quella della corona
esterna. Chiedersi quanto valga la singola feature numero~87 --- il sensore di
nord-est di cinque passi fa --- non ha una risposta sensata, perch\'e essa \`e quasi
determinata da una dozzina di altre. Il problema non \`e solo di leggibilit\`a: in
presenza di feature perfettamente correlate il valore di Shapley ripartisce il merito
fra di esse in parti uguali, diluendo l'attribuzione di un blocco importante su decine
di componenti. KernelSHAP aggiunge a ci\`o un'ipotesi di indipendenza fra le feature nel
campionamento delle coalizioni, palesemente violata qui, con distorsioni
documentate~\cite{aas2021}. La risposta corretta non \`e abbandonare Shapley ma
cambiarne i giocatori: \`e la costruzione nota come
\emph{groupShapley}~\cite{groupshapley}, in cui i giocatori sono gruppi di feature
semanticamente coerenti. I gruppi adottati sono definiti nella
Sezione~\ref{sec:xai_gruppi} e sono, non a caso, le stesse unit\`a di cui si parla nel
Capitolo~\ref{cap:rl}: la seconda corona, il vento, la memoria degli spostamenti.

\paragraph{Esattezza.} Con dodici gruppi le coalizioni sono $2^{12}=4096$: un numero
enumerabile in pochi secondi. Gli stimatori servono quando l'enumerazione \`e
impraticabile; qui non lo \`e, e i valori di Shapley si calcolano in forma
\emph{esatta}. Si evita cos\`i sia l'errore di approssimazione sia la scelta arbitraria
del numero di campioni, e l'assioma di efficienza diventa un test numerico di
correttezza del calcolo anzich\'e una propriet\`a asserita
(Sezione~\ref{sec:xai_shapley}).

\paragraph{La possibilit\`a di intervenire.} \`E la ragione pi\`u profonda. Nel
supervised learning i metodi post-hoc sono necessari perch\'e il mondo non \`e
rieseguibile: il modello \`e dato, i dati sono dati, e l'importanza di una feature pu\`o
solo essere \emph{stimata} da un surrogato --- da cui la \emph{model
reliance}~\cite{model_reliance} e le sue varianti a permutazione. Nel reinforcement
learning con simulatore questa limitazione non sussiste. Bloccare un gruppo di feature
e rieseguire la politica misurando il success rate effettivo non \`e un surrogato
dell'importanza: \`e un \emph{intervento}, cio\`e evidenza diretta. Inoltre la politica,
a differenza di un classificatore, non \`e un fine ma un mezzo per un ritorno
misurabile, il che rende disponibile una metrica di importanza definita sul compito e
non sulla funzione di perdita --- uno status che l'ablation ha nell'explainability per
RL e non nel supervised~\cite{xrl_survey}. Rinunciarvi per affidarsi soltanto a metodi
osservazionali significherebbe lasciare inutilizzata la prova migliore a disposizione.

\paragraph{Perch\'e non LIME.} LIME \`e escluso come strumento principale per ragioni
indipendenti. \`E \emph{locale} per costruzione: spiega una singola decisione, mentre le
domande di questo capitolo sono globali (il vento viene usato? la seconda corona
serve?); rispondervi richiederebbe di aggregare centinaia di spiegazioni locali,
ricostruendo in modo deteriore un metodo globale. Le sue perturbazioni campionano
inoltre le feature indipendentemente, il che con blocchi collineari produce punti
lontani dalla variet\`a dei dati --- la storia dei sensori \`e quasi deterministica dati
i sensori correnti. Infine le spiegazioni LIME sono note per essere instabili rispetto
alla larghezza del kernel e al seed di campionamento~\cite{robustness_interp}: fondare
su di esse un'affermazione del tipo ``la feature $X$ conta'' \`e metodologicamente
fragile. Resta a LIME una nicchia legittima, che questo capitolo non sfrutta ma che
segnaliamo fra gli sviluppi possibili: lo studio di caso di una singola decisione
critica, dove la domanda \`e effettivamente locale.

\paragraph{Il criterio adottato.} Ne segue l'impianto del capitolo: valori di Shapley
\emph{esatti} sui gruppi per stabilire che cosa la politica guarda, ablation
\emph{causale} sugli stessi gruppi per stabilire che cosa serve al compito, e infine
sonde comportamentali per caratterizzare la strategia appresa. I tre livelli sono
complementari e vengono confrontati nella Sezione~\ref{sec:xai_sintesi}.

\subsection{Due accortezze tecniche}
\label{sec:xai_accortezze}

Due dettagli dell'implementazione hanno conseguenze sostanziali sui risultati e vanno
dichiarati.

\textbf{Logit pre-maschera.} MaskablePPO~\cite{maskable_ppo} pone a $-\infty$ i logit
delle azioni che porterebbero l'agente a terra o fuori dominio
(Sezione~\ref{sec:osservazione}). Attribuzioni e sonde comportamentali sono calcolate
sui logit \emph{prima} del mascheramento: diversamente si misurerebbe la geometria
della costa invece della politica. L'ablation, al contrario, gira con il mascheramento
attivo, perch\'e in quel caso l'oggetto di interesse \`e il compito reale, di cui il
mascheramento \`e parte.

\textbf{Spazio normalizzato.} La rete non vede le grandezze fisiche ma la loro versione
normalizzata da \texttt{VecNormalize}, che sottrae la media e divide per la deviazione
standard stimate in addestramento. Ci\`o ha due implicazioni favorevoli. Primo, tutte le
feature hanno varianza unitaria, quindi le magnitudini di attribuzione sono
confrontabili fra blocchi senza ulteriori riscalature. Secondo, il vettore nullo nello
spazio normalizzato corrisponde esattamente alla \emph{media di addestramento}: \`e
quindi una baseline naturale, con la semantica di ``sensore che riporta il valore
tipico'', e la stessa baseline serve sia per le coalizioni vuote dello Shapley sia per
l'ablation. Fa eccezione la lettura dei sensori grezzi nelle sonde comportamentali,
dove la normalizzazione per-feature altererebbe l'ordinamento e va quindi invertita.

\section{I gruppi semantici dell'osservazione}
\label{sec:xai_gruppi}

L'analisi \`e condotta sul modello a \textbf{doppia corona con $v_{\max} = 2\,$m/s}
(Sezione~\ref{sec:rl_dual}), che ha osservazione a 196 dimensioni e spazio d'azione
$\text{Discrete}(8\times5)$. Le 196 feature sono raggruppate nei dodici blocchi della
Tabella~\ref{tab:xai_gruppi}, che ricalcano la struttura dell'osservazione
(Tabella~\ref{tab:obs}). La valutazione usa il protocollo held-out del
Capitolo~\ref{cap:rl}: 26 sorgenti mai viste, 4 scenari di vento, 3 chunk temporali.

\begin{table}[H]
\centering
\renewcommand{\arraystretch}{1.2}
\small
\begin{tabular}{L{5.6cm} C{2.0cm} C{1.6cm}}
\toprule
\rowcolor{headerbg}
\textcolor{white}{\textbf{Gruppo}} &
\textcolor{white}{\textbf{Indici}} &
\textcolor{white}{\textbf{Dim.}} \\
\midrule
Concentrazione (istante)   & [0]        & 1  \\
\rowcolor{rowalt}
Concentrazione (storia)    & [1:10]     & 9  \\
Spostamenti (storia)       & [10:28]    & 18 \\
\rowcolor{rowalt}
Corona 20\,m (istante)     & [28:36]    & 8  \\
Corona 50\,m (istante)     & [36:44]    & 8  \\
\rowcolor{rowalt}
Corona 20\,m (storia)      & [44:116]   & 72 \\
Corona 50\,m (storia)      & [116:188]  & 72 \\
\rowcolor{rowalt}
Vento $(u,v)$              & [188:190]  & 2  \\
Corrente $(u,v)$           & [190:192]  & 2  \\
\rowcolor{rowalt}
Vettore verso max conc.    & [192:194]  & 2  \\
Valore max conc.           & [194]      & 1  \\
\rowcolor{rowalt}
Et\`a ultimo contatto      & [195]      & 1  \\
\midrule
\textbf{Totale}            &            & \textbf{196} \\
\bottomrule
\end{tabular}
\caption{I dodici gruppi semantici in cui \`e partizionata l'osservazione a doppia
         corona. Sono i giocatori dello Shapley (Sezione~\ref{sec:xai_shapley}) e le
         unit\`a dell'ablation (Sezione~\ref{sec:xai_ablation}).}
\label{tab:xai_gruppi}
\end{table}

\section{Ablation causale di gruppo}
\label{sec:xai_ablation}

\paragraph{Protocollo.} Per ogni gruppo si azzerano le sue componenti nello spazio
normalizzato --- equivalente, come detto, a sostituirle con la media di addestramento ---
e si riesegue la politica misurando success rate, step al successo e distanza finale. Il
disegno \`e \textbf{appaiato}: per ogni spawn si valutano la baseline e tutte e dodici le
ablazioni a partire dallo stesso seme, quindi dallo stesso punto di partenza e dallo
stesso campo. Senza questo accorgimento la varianza fra spawn --- la distanza iniziale
varia da 300 a oltre 2000\,m --- sommergerebbe differenze di pochi punti. Sono stati
valutati \textbf{624 spawn appaiati} $\times$ 13 condizioni, per un totale di circa 8100
episodi; gli intervalli di confidenza al 95\% sono ottenuti per bootstrap
ricampionando \emph{gli spawn} (non gli episodi), cos\`i da preservare l'appaiamento.
L'ablazione di un gruppo si riduce ad azzerarne gli indici nell'osservazione gi\`a
normalizzata, ad ogni passo dell'episodio:

\begin{lstlisting}
def run_episode_ablated(model, vec_env, ablate_idx, seed):
    obs = vec_env.reset()                  # obs normalizzata (VecNormalize)
    obs[:, ablate_idx] = 0.0               # azzera il gruppo -> media di training
    done = False
    while not done:
        action, _ = model.predict(obs, deterministic=True,
                                  action_masks=vec_env.env_method("action_masks")[0])
        obs, _, done, info = vec_env.step(action)
        obs[:, ablate_idx] = 0.0           # mantiene il blocco azzerato ad ogni step
    return info
\end{lstlisting}

\begin{figure}[H]
  \centering
  \includegraphics[width=\linewidth]{fig_ablation_clean_dualcorona.png}
  \caption{Ablation causale di gruppo sul modello a doppia corona $v_{\max}{=}2$: 624
           spawn appaiati, IC 95\% bootstrap. A sinistra il costo in efficienza (step al
           successo), a destra quello in riuscita (success rate). L'ordinamento segue
           l'efficienza. Il colore indica la \emph{classe di evidenza}: rosso se
           l'intervallo esclude lo zero, grigio se lo contiene.}
  \label{fig:xai_ablation}
\end{figure}

\paragraph{Risultati.} La baseline raggiunge il $99{,}8\%$ di success rate in $47{,}6$
step medi --- valori misurati sui 624 spawn appaiati di questo protocollo, e quindi
lievemente diversi dal $99{,}7\%$ / $49{,}8$ step riportati per lo stesso modello nella
Tabella~\ref{tab:dual_vmax} (valutazione held-out a 1560 episodi). Il risultato dominante \`e la \textbf{storia della corona a 20\,m}: rimuoverla
costa $-26{,}8$ punti di success rate (IC $[-30{,}3, -23{,}4]$) e $+77{,}7\%$ di step. \`E
di gran lunga il blocco pi\`u portante dell'intera osservazione, e il fatto che a pesare
sia la \emph{storia} --- non la lettura istantanea, che costa $-6{,}2$ punti --- dice che
la politica integra il campo chimico nel tempo anzich\'e reagire all'istante. Seguono la
memoria degli \textbf{spostamenti} ($-12{,}8$ punti), cio\`e il senso di dove si sta
andando, e il \textbf{vettore verso la massima concentrazione vista} ($-5{,}0$).

All'estremo opposto, tre gruppi risultano \textbf{indistinguibili dalla baseline} su
entrambe le metriche: la concentrazione istantanea, la sua storia e la corona a 50\,m
istantanea hanno intervalli di confidenza che contengono lo zero. Il caso della
\textbf{concentrazione istantanea} \`e il pi\`u netto: azzerarla non cambia nulla
($-0{,}7\%$ di step, IC $[-5{,}5, +4{,}5]$; success rate invariato). La lettura al centro
della formazione \`e ridondante --- gli otto sensori direzionali la circondano a
$20\,\text{m}$ e ne contengono gi\`a l'informazione --- e la politica, di fatto, la
ignora.

\paragraph{La seconda corona.} Il verdetto sulla corona a 50\,m \`e sfumato e merita
attenzione, perch\'e incrocia il risultato della Sezione~\ref{sec:rl_dual}. La sua lettura
istantanea non ha effetto misurabile ($+5{,}3\%$ di step, IC $[-0{,}4, +12{,}4]$), mentre
la sua \emph{storia} ha un effetto piccolo ma reale sull'efficienza ($+13{,}9\%$, IC
$[+7{,}0, +21{,}9]$) e nullo sulla riuscita ($-0{,}3$ punti). Il quadro \`e coerente con
il confronto per riaddestramento del Capitolo~\ref{cap:rl}, dove la doppia corona a
$v_{\max}{=}2$ risultava $28{,}3\%$ pi\`u veloce a parit\`a di success rate: i due
esperimenti --- rimuovere il sensore da un modello \emph{congelato} e riaddestrare senza
di esso --- convergono sullo stesso ordine di grandezza e sulla stessa conclusione
qualitativa, cio\`e che la seconda corona compra \textbf{velocit\`a, non affidabilit\`a}.
La convergenza di due protocolli indipendenti \`e un riscontro non banale.

\section{Valori di Shapley esatti}
\label{sec:xai_shapley}

Con dodici giocatori le $2^{12} = 4096$ coalizioni sono state enumerate integralmente su
un buffer di 1500 stati raccolti dai rollout della politica piena. Il payoff \`e
$v(S) = \pi(a^*\,|\,x_S)$, con $a^*$ l'azione scelta a osservazione completa. La media di
$\pi(a^*)$ passa da $0{,}201$ alla baseline a $0{,}707$ a osservazione piena: i dodici
gruppi si ripartiscono quindi $0{,}506$ di massa di probabilit\`a. L'assioma di
efficienza \`e verificato numericamente con un errore massimo di
$3{,}2 \times 10^{-8}$, che conferma l'esattezza del calcolo.

La Figura~\ref{fig:xai_shapley} riporta, per ciascuno dei dodici gruppi, la magnitudine media
del valore di Shapley $|\phi_i|$: la lunghezza di ogni barra misura quanto quel blocco
contribuisce, in media, alla probabilit\`a che la politica assegna all'azione scelta. \`E
quindi una graduatoria di \emph{ci\`o che la rete guarda}, ossia una lettura
\emph{osservazionale} della funzione appresa, distinta dalla \emph{necessit\`a} causale
misurata dall'ablazione (Sezione~\ref{sec:xai_ablation}); come si vedr\`a, le due letture
divergono, e il confronto \`e ripreso nella Sezione~\ref{sec:xai_sintesi}.

\begin{figure}[H]
  \centering
  \includegraphics[width=0.94\linewidth]{fig_shapley_clean_dualcorona.png}
  \caption{Valori di Shapley esatti sui dodici gruppi: media di $|\phi_i|$ su 1500 stati,
           $4096$ coalizioni enumerate. La magnitudine \`e sull'asse, non nel colore.}
  \label{fig:xai_shapley}
\end{figure}

La graduatoria di ci\`o che la rete \emph{guarda} \`e guidata a pari merito dalla
\textbf{corrente} ($|\phi| = 0{,}236$) e dalla \textbf{memoria degli spostamenti}
($0{,}236$), seguite dalla storia della corona a 20\,m ($0{,}183$). In coda stanno l'et\`a
dell'ultimo contatto ($0{,}023$), la corona a 50\,m istantanea ($0{,}031$) e la
concentrazione istantanea ($0{,}038$).

\section{Analisi comportamentale}
\label{sec:xai_comportamento}

\paragraph{Non \`e risalita di gradiente.} La domanda pi\`u naturale su una politica di
source seeking \`e se si riduca a seguire il sensore pi\`u intenso. Per rispondere si
confronta, su 13\,738 stati con contatto col plume, la direzione del sensore massimo
della corona a 20\,m con la direzione scelta. La lettura richiede un'accortezza: gli spawn
stanno sul plume, che \`e avvettato \emph{a valle} della sorgente, quindi la sorgente cade
quasi sempre in una direzione sistematica e il marginale $P(\text{azione})$ \`e
fortemente sbilanciato. Confrontare la diagonale con $1/8$ sarebbe quindi fuorviante: va
confrontata con il valore \emph{atteso sotto indipendenza}, che incorpora quel bias.

\begin{figure}[H]
  \centering
  \includegraphics[width=\linewidth]{fig_sensore_azione.png}
  \caption{Sensore massimo (corona 20\,m) contro direzione scelta. A sinistra la
           condizionale $P(a|s)$, dominata dal bias direzionale degli scenari; a destra il
           \emph{lift} $P(a|s)/P(a)$, che divide per il marginale e isola il solo
           contributo del sensore ($1 = $ nessun effetto).}
  \label{fig:xai_sensore}
\end{figure}

La massa sulla diagonale \`e $0{,}068$ contro un atteso di $0{,}064$: un \textbf{lift di
appena $1{,}06\times$}. La politica muove verso il sensore pi\`u intenso solo il $6\%$
delle volte in pi\`u di quanto predica il suo stesso bias direzionale. \textbf{La politica
non fa risalita di gradiente locale}: pur essendo la corona a 20\,m il blocco pi\`u
portante dell'osservazione (Sezione~\ref{sec:xai_ablation}), non la usa per scegliere la
direzione istantanea di massima concentrazione. I due risultati non sono in
contraddizione, anzi si spiegano a vicenda: ci\`o che conta non \`e la corona ma la
\emph{sua storia}, cio\`e l'integrazione temporale del campo, esattamente la strategia che
il Capitolo~\ref{cap:rl} attribuiva all'RL per spiegare il divario con l'FCM --- il quale,
per costruzione, la risalita di gradiente locale la fa.

\paragraph{Il vento \`e letto ma non seguito.} Ruotando il vento a $360^\circ$ e tenendo
fermo il resto dell'osservazione, la direzione preferita dalla politica insegue il
controvento con un allineamento di appena $0{,}109$ (dove $1$ sarebbe anemotassi perfetta
e $0$ indifferenza). La politica, cio\`e, \textbf{non fa anemotassi}. Il dato \`e coerente
con l'addestramento: questo modello discende dalla catena \texttt{minimal}, in cui i
reward di allineamento a vento e corrente sono \emph{disattivati}
(Sezione~\ref{sec:rl_reward}) proprio perch\'e controproducenti. Il risultato vale quindi
anche come validazione del metodo: l'analisi ritrova, senza saperlo, una propriet\`a nota
del setup di addestramento.

Va segnalato il limite della sonda: ruotare il vento tenendo fermo il campo di
concentrazione produce osservazioni fuori dalla variet\`a dei dati, perch\'e vento e forma
del plume sono fisicamente accoppiati. La response curve risponde quindi alla domanda
``l'uscita della politica dipende dall'ingresso vento?'', non a ``cosa accadrebbe in un
mondo con vento diverso?''.

\paragraph{Il critic stima la distanza.} La funzione valore $V(s)$ correla negativamente
con la distanza dalla sorgente ($r = -0{,}33$ su tutti gli scenari): il critic ha appreso
una nozione di prossimit\`a senza che gliela si sia mai fornita esplicitamente. La
correlazione \`e moderata perch\'e aggregata su 26 sorgenti e 12 combinazioni
vento--chunk, ciascuna con il proprio livello di difficolt\`a: $V(s)$ stima il ritorno
atteso, che dipende anche dallo scenario, non solo dalla distanza.

\begin{figure}[H]
  \centering
  \includegraphics[width=\linewidth]{fig_valore_entropia.png}
  \caption{A sinistra $V(s)$ stimato dal critic in funzione della distanza dalla sorgente
           (mediana e intervallo interquartile per bin); a destra l'entropia della
           politica, calcolata sui logit pre-maschera.}
  \label{fig:xai_valore}
\end{figure}

\section{Il valore della formazione: ablazione per riaddestramento}
\label{sec:xai_noform}

Le ablazioni fin qui condotte operano su un modello \textbf{congelato}: si azzera un gruppo
di feature in inferenza e si misura il danno. Questo protocollo, per\`o, \emph{sovrastima}
l'importanza di una feature, perch\'e l'agente non pu\`o adattarsi alla sua assenza. Per
misurare il valore \emph{irriducibile} della formazione --- l'insieme dei sensori
direzionali che distingue la nostra squadra master--slave da un singolo veicolo --- si
adotta un protocollo pi\`u equo: \textbf{rimuovere la formazione e riaddestrare}.

\paragraph{Protocollo.} A partire dal modello a singola corona con $v_{\max}=2\,\text{m/s}$
si esegue un fine-tuning ($2$\,M step) in cui i sensori direzionali della formazione ---
correnti e storici --- sono \textbf{azzerati} ad ogni passo (\emph{masking} della
formazione). L'architettura non cambia (osservazione a 116 dimensioni, spazio d'azione
$\text{Discrete}(8\times5)$): l'agente riparte dai pesi del modello con formazione e impara
una nuova strategia con la sola informazione che gli resta --- la concentrazione nella
\emph{propria} posizione, le memorie di concentrazione e di spostamento, il vento e la
corrente. Tutto il resto (velocit\`a, protocollo held-out a 1560 episodi, dataset)
\`e identico, cos\`i che il calo di prestazioni sia attribuibile alla \emph{sola} rimozione
della formazione.

\paragraph{Costo della rimozione.} La Tabella~\ref{tab:noform_cmp} confronta l'agente
singolo con il modello a formazione di pari velocit\`a.

\begin{table}[H]
\centering
\renewcommand{\arraystretch}{1.3}
\small
\begin{tabular}{L{3.4cm} C{2.8cm} C{2.8cm} C{2.4cm}}
\toprule
\rowcolor{headerbg}
\textcolor{white}{\textbf{Metrica}} &
\textcolor{white}{\textbf{Con formazione}} &
\textcolor{white}{\textbf{Senza formazione}} &
\textcolor{white}{\textbf{$\Delta$}} \\
\midrule
\rowcolor{rowalt} \textbf{SR globale} & 99{,}7\% & \textbf{93{,}3\%} & $-6{,}4$ \\
Step medi (succ.) & 69{,}5 & 100{,}6 & $+44{,}7\%$ \\
\midrule
V0 & 100{,}0\% & 97{,}2\% & $-2{,}8$ \\
\rowcolor{rowalt} V1 & 99{,}5\% & 97{,}9\% & $-1{,}6$ \\
V2 & 99{,}2\% & 86{,}4\% & $\mathbf{-12{,}8}$ \\
\rowcolor{rowalt} V3 & 100{,}0\% & 91{,}8\% & $-8{,}2$ \\
Q1/4 & 100{,}0\% & 97{,}9\% & $-2{,}1$ \\
\rowcolor{rowalt} Q1/2 & 99{,}0\% & 90{,}0\% & $\mathbf{-9{,}0}$ \\
Q3/4 & 100{,}0\% & 92{,}1\% & $-7{,}9$ \\
\bottomrule
\end{tabular}
\caption{Modello con vs.\ senza formazione a $v_{\max}{=}2\,$m/s (1560 episodi ciascuno). Il
         $\Delta$ \`e il costo \emph{netto} della rimozione della formazione, dopo che
         l'agente ha avuto modo di adattarsi.}
\label{tab:noform_cmp}
\end{table}

La formazione vale \textbf{$6{,}4$ punti} di success rate globale e circa il $45\%$ in
velocit\`a. Il costo non \`e uniforme: \`e trascurabile negli scenari facili (V0/V1, Q1/4,
entro $-3$ punti) e diventa marcato dove il plume \`e avvettato e diffuso --- \textbf{V2}
($-12{,}8$) e \textbf{Q1/2} ($-9{,}0$) --- ossia dove il gradiente \emph{spaziale}
istantaneo delle otto letture simultanee \`e pi\`u prezioso. Il calo resta comunque
\textbf{moderato}: l'agente singolo riaddestrato recupera gran parte delle prestazioni, e
soprattutto rimane molto sopra il metodo del gradiente --- l'FCM adattivo a pari velocit\`a
($lr=20\,$m) si ferma al $58{,}2\%$, circa \textbf{35 punti} sotto l'agente singolo (93{,}3\%).
\`E un fatto notevole: un agente RL che vede solo la propria concentrazione batte nettamente
un metodo reattivo che dispone dell'intera formazione. Ci\`o che separa il successo dal
fallimento non \`e il sensore, ma la strategia appresa.

\paragraph{Come compensa l'agente singolo.} L'explainability del modello senza formazione
(stesso protocollo: ablazione di gruppo e Shapley esatti) rivela uno \textbf{spostamento del
baricentro informativo dal sensing spaziale a quello temporale}
(Tabella~\ref{tab:noform_shift}, Figura~\ref{fig:noform_xai}). Privato della corona,
l'agente si regge in modo schiacciante sulla \textbf{memoria degli spostamenti} (ablazione
$-55{,}6$ punti, Shapley $|\phi|=0{,}27$): correla i propri movimenti con le variazioni
della concentrazione che misura sotto di s\'e, ricostruendo nel \emph{tempo} il gradiente
che la formazione forniva nello \emph{spazio}. Seguono l'\emph{homing} verso la massima
concentrazione vista ($-25{,}8$), la corrente ($-12{,}8$) e la storia della concentrazione
propria ($-11{,}2$) --- proprio la feature che il modello \emph{con} formazione ignorava
($\approx 0$).

Come per il modello con formazione (Sezione~\ref{sec:xai_sintesi}), anche qui
\emph{attribuzione} e \emph{necessit\`a} divergono, e nello stesso modo istruttivo. Il
\textbf{vettore verso il massimo di concentrazione} \`e solo quarto per attribuzione
($|\phi| = 0{,}07$) ma \textbf{secondo per impatto causale} ($-25{,}8$): l'homing di memoria
\`e insostituibile per il compito pur essendo poco ``guardato'' dalla rete. Specularmente, la
\textbf{corrente} \`e seconda per attribuzione ($|\phi| = 0{,}12$) ma pesa meno all'ablazione
($-12{,}8$), confermando --- ora anche in assenza della corona --- che ci\`o che la rete
osserva non coincide con ci\`o che le \`e strettamente necessario.

\begin{table}[H]
\centering
\renewcommand{\arraystretch}{1.3}
\small
\begin{tabular}{L{5.4cm} C{2.9cm} C{2.9cm}}
\toprule
\rowcolor{headerbg}
\textcolor{white}{\textbf{Gruppo ($\Delta$SR)}} &
\textcolor{white}{\textbf{Con formazione}} &
\textcolor{white}{\textbf{Senza formazione}} \\
\midrule
Corona 20\,m (storia)              & $\mathbf{-26{,}8}$ & mascherata \\
\rowcolor{rowalt} Spostamenti (storia) & $-12{,}8$ & $\mathbf{-55{,}6}$ \\
Corona 20\,m (istante)             & $-6{,}2$  & mascherata \\
\rowcolor{rowalt} Vettore verso max conc.\ (homing) & $-5{,}0$ & $-25{,}8$ \\
Corrente $(u,v)$                   & $-3{,}0$  & $-12{,}8$ \\
\rowcolor{rowalt} Vento $(u,v)$    & $-2{,}2$  & $-1{,}4$ \\
Concentrazione (storia)            & $-0{,}5$  & $-11{,}2$ \\
\rowcolor{rowalt} Corona 50\,m (storia) & $-0{,}3$ & --- \\
Concentrazione (istante)           & $0{,}0$   & $-3{,}2$ \\
\rowcolor{rowalt} Corona 50\,m (istante) & $0{,}0$ & --- \\
Valore max conc.                   & $+0{,}2$  & $-7{,}1$ \\
\rowcolor{rowalt} Et\`a ultimo contatto & $+0{,}2$ & $-0{,}5$ \\
\bottomrule
\end{tabular}
\caption{Spostamento del baricentro informativo: costo in success rate ($\Delta$SR,
         punti percentuali) di \emph{ogni} gruppo, per il modello \emph{con} formazione
         (doppia corona, ablazioni della Sezione~\ref{sec:xai_ablation}) e per l'agente
         \emph{senza} formazione (singola corona, formazione mascherata). Con la formazione la
         politica si regge sul gradiente \emph{spaziale} (storia della corona); rimossa, il
         peso si trasferisce sul sensing \emph{temporale} --- spostamenti, homing, storia e
         valore di concentrazione, tutti nettamente pi\`u costosi. ``mascherata'': corona
         azzerata, effetto spurio escluso (Figura~\ref{fig:noform_xai}); ``---'': gruppo
         assente, poich\'e l'agente singolo ha una sola corona.}
\label{tab:noform_shift}
\end{table}

\begin{figure}[H]
  \centering
  \begin{subfigure}{\linewidth}
    \centering
    \includegraphics[width=\linewidth]{fig_ablation_clean_noformation.png}
    \caption{Ablazione di gruppo.}
  \end{subfigure}

  \vspace{5mm}
  \begin{subfigure}{\linewidth}
    \centering
    \includegraphics[width=0.9\linewidth]{fig_shapley_clean_noformation.png}
    \caption{Valori di Shapley esatti.}
  \end{subfigure}
  \caption{Explainability dell'agente singolo (senza formazione), sui soli gruppi
           \emph{genuini}: la memoria degli spostamenti \`e la feature portante, seguita
           dall'homing verso il massimo di concentrazione e dalla corrente. I gruppi della
           corona sono \textbf{esclusi} perch\'e mascherati a zero: \texttt{VecNormalize} li
           rende una costante non nulla il cui azzeramento perturberebbe il bias appreso,
           producendo un effetto spurio privo di contenuto informativo.}
  \label{fig:noform_xai}
\end{figure}

Il quadro chiude il cerchio dell'analisi: la rete non ``perde un pezzo'' quando le si toglie
la formazione, ma \textbf{ricompone} la stima della direzione con gli strumenti che le
restano. La formazione \`e un vantaggio reale ma \emph{sostituibile in parte} da una
strategia temporale appresa; \`e negli scenari a plume diffuso, dove il segnale temporale
degenera pi\`u in fretta, che il suo contributo spaziale resta insostituibile.

\section{Sintesi: che cosa la rete guarda, che cosa le serve}
\label{sec:xai_sintesi}

Il confronto fra le due graduatorie \`e la parte pi\`u istruttiva del capitolo, perch\'e
\textbf{divergono}.

\begin{table}[H]
\centering
\renewcommand{\arraystretch}{1.3}
\small
\begin{tabular}{L{4.1cm} C{2.2cm} C{1.9cm} L{3.9cm}}
\toprule
\rowcolor{headerbg}
\textcolor{white}{\textbf{Gruppo}} &
\textcolor{white}{\textbf{Shapley $|\phi|$}} &
\textcolor{white}{\textbf{$\Delta$SR}} &
\textcolor{white}{\textbf{Lettura}} \\
\midrule
Corona 20\,m (storia)          & 0{,}183 (3\textsuperscript{o})  & $\mathbf{-26{,}8}$ & indispensabile \\
\rowcolor{rowalt}
Spostamenti (storia)           & 0{,}236 (1\textsuperscript{o})  & $-12{,}8$ & indispensabile \\
Corona 20\,m (istante)         & 0{,}104 (4\textsuperscript{o})  & $-6{,}2$  & necessaria \\
\rowcolor{rowalt}
Vettore verso max conc.\ (homing) & 0{,}044 (8\textsuperscript{o}) & $-5{,}0$ & necessaria, poco guardata \\
Corrente $(u,v)$               & \textbf{0{,}236} (1\textsuperscript{o}) & $-3{,}0$ & guardata, poco necessaria \\
\rowcolor{rowalt}
Vento $(u,v)$                  & 0{,}075 (6\textsuperscript{o})  & $-2{,}2$  & guardata, poco necessaria \\
Concentrazione (storia)        & 0{,}082 (5\textsuperscript{o})  & $-0{,}5$  & guardata, poco necessaria \\
\rowcolor{rowalt}
Corona 50\,m (storia)          & 0{,}058 (7\textsuperscript{o})  & $-0{,}3$  & ridondante \\
Concentrazione (istante)       & 0{,}038 (10\textsuperscript{o}) & $0{,}0$   & ridondante \\
\rowcolor{rowalt}
Corona 50\,m (istante)         & 0{,}031 (11\textsuperscript{o}) & $0{,}0$   & ridondante \\
Valore max conc.               & 0{,}044 (8\textsuperscript{o})  & $+0{,}2$  & ridondante \\
\rowcolor{rowalt}
Et\`a ultimo contatto          & 0{,}023 (12\textsuperscript{o}) & $+0{,}2$  & ridondante \\
\bottomrule
\end{tabular}
\caption{Attribuzione (Shapley $|\phi|$: \emph{che cosa la rete guarda}) contro necessit\`a
         causale (ablation $\Delta$SR in punti percentuali: \emph{che cosa serve al compito}),
         per tutti i dodici gruppi, ordinati per necessit\`a decrescente. Le due graduatorie
         \textbf{divergono}: la \emph{corrente} \`e la pi\`u guardata (Shapley 1\textsuperscript{o})
         ma poco necessaria ($-3{,}0$), mentre la \emph{storia della corona} \`e solo terza per
         attribuzione ma di gran lunga la pi\`u necessaria ($-26{,}8$); specularmente l'\emph{homing}
         \`e poco guardato (8\textsuperscript{o}) ma utile ($-5{,}0$).}
\label{tab:xai_sintesi}
\end{table}

Il caso pi\`u istruttivo \`e la \textbf{corrente marina}: \`e il gruppo con
l'attribuzione pi\`u alta in assoluto ($|\phi| = 0{,}236$, primo a pari merito), eppure
rimuoverla costa appena $3$ punti di success rate. La rete la legge molto, ma il compito
ne ha poco bisogno --- segno che l'informazione che la corrente veicola \`e in larga parte
\emph{recuperabile} dagli altri blocchi, in primo luogo dalla storia della corona, che del
trasporto porta gi\`a la traccia. Simmetricamente, la \textbf{storia della corona a
20\,m} \`e solo terza per attribuzione ma prima, e di gran lunga, per impatto causale: \`e
l'unico blocco davvero non sostituibile.

\`E precisamente la divergenza che giustifica l'impianto metodologico della
Sezione~\ref{sec:xai_rassegna}. Un capitolo fondato sui soli metodi post-hoc avrebbe
concluso che la corrente \`e la grandezza pi\`u importante per la politica; l'intervento
mostra che, se la si toglie, la politica se la cava quasi altrettanto bene. Le due
domande --- che cosa la rete guarda, che cosa serve al compito --- hanno risposte diverse,
e solo la seconda \`e quella che conta per decidere quali sensori montare a bordo.

Ne emerge un ritratto coerente della politica appresa. Essa \textbf{non} \`e un
risalitore di gradiente: la direzione istantanea di massima concentrazione non guida la
scelta (lift $1{,}06\times$), e la concentrazione al centro della formazione \`e
addirittura ignorata. Ci\`o su cui si regge \`e l'\textbf{integrazione temporale} del
campo chimico locale --- la storia della corona a 20\,m, il blocco pi\`u portante, unita
alla memoria di dove si \`e andati --- affiancata da una memoria esplicita del punto
migliore incontrato. La seconda corona a 50\,m serve a puntare meglio le falcate lunghe,
comprando velocit\`a e non affidabilit\`a, in accordo con quanto il
Capitolo~\ref{cap:rl} aveva stabilito per riaddestramento. Vento e corrente vengono letti
ma non usati come bussola. \`E, in sintesi, esattamente la strategia che il
Capitolo~\ref{cap:rl} ipotizzava per spiegare il divario con il metodo del gradiente ---
con la differenza che qui \`e misurata anzich\'e supposta.
